\section{Introduction}
\label{sec:intro}

Detecting visual anomalies in industrial manufacturing is a critical yet challenging task in computer vision~\cite{Bergmann2019MVTec}.
In production lines, even subtle surface defects such as scratches, dents, or contaminations can lead to product failures and significant financial losses.
Automated visual inspection systems are therefore essential for maintaining quality standards and minimizing human oversight costs.

The primary obstacle in industrial anomaly detection (AD) is the scarcity of labeled anomaly data.
Anomalies are inherently rare, diverse, and unpredictable, making it impractical to collect comprehensive anomaly datasets.
Consequently, most practical AD systems adopt a one-class classification paradigm: only normal (defect-free) samples are available during training, and any deviation from the learned normality model is flagged as anomalous at test time.
Within this paradigm, the \emph{few-shot} setting---where only $k \in \{1, 2, 4\}$ normal images per category are available---represents the most realistic and challenging scenario~\cite{Huang2022WinCLIP,Li2024PromptAD}.

Recent advances in few-shot anomaly detection have produced impressive results, driven largely by the emergence of powerful vision foundation models such as DINOv2~\cite{Oquab2023DINOv2} and CLIP~\cite{Radford2021CLIP}.
The recently proposed SubspaceAD~\cite{Lendering2025SubspaceAD} exemplifies a minimalist yet highly effective approach: it extracts dense patch features using a frozen DINOv2 backbone, fits a Principal Component Analysis (PCA) model to the normal features, and detects anomalies via reconstruction residuals with respect to the learned linear subspace.
Despite its simplicity, SubspaceAD achieves state-of-the-art results across multiple few-shot settings, outperforming substantially more complex memory-bank-based, reconstruction-based, and vision-language model approaches.

The success of SubspaceAD raises an intriguing question: \emph{Can we further improve anomaly detection by replacing the linear PCA subspace with a more expressive nonlinear manifold model, while preserving the minimalist, efficient philosophy of the original framework?}

In this paper, we answer this question affirmatively by introducing \textbf{MSAE-AD}, a novel anomaly detection framework that substitutes PCA with a Matryoshka Sparse Autoencoder (MSAE)~\cite{Zaigrajew2025MSAE}.
MSAE is a recently proposed hierarchical sparse autoencoder architecture that learns representations at multiple granularities simultaneously, establishing a state-of-the-art Pareto frontier between reconstruction quality and sparsity.
Our key insight is that MSAE's hierarchical sparse representations can model the normal feature manifold more accurately than a single linear subspace, leading to more precise anomaly discrimination.

Specifically, our contributions are as follows:
\begin{itemize}
    \item We propose MSAE-AD, a novel few-shot anomaly detection framework that leverages MSAE for multi-scale, nonlinear modeling of normal feature distributions, replacing the PCA module in SubspaceAD while retaining the same DINOv2 feature extraction backbone.
    \item We investigate two MSAE training variants---uniform weighting (UW) and reverse weighting (RW)---and analyze their respective advantages for anomaly detection tasks.
    \item We conduct comprehensive experiments on the MVTec-AD and VisA benchmarks, demonstrating that MSAE-AD achieves strong detection and localization performance across all 15 MVTec-AD categories and 12 VisA categories.
    \item We provide detailed per-category analysis and discuss the complementary strengths of linear (PCA) and nonlinear (MSAE) approaches to normal manifold modeling.
\end{itemize}

Our experimental results show that MSAE-AD with RW weighting achieves average image-level AUROC scores of 98.72\% on MVTec-AD and 95.58\% on VisA, with pixel-level AUROC of 97.71\% and 98.46\% respectively.
The UW variant achieves slightly lower but highly competitive image-level scores (98.07\% on MVTec-AD, 95.08\% on VisA) while offering better pixel-level PRO metrics in certain settings.
These results establish MSAE-AD as an effective and principled approach that extends the minimalist SubspaceAD paradigm through hierarchical sparse modeling.
